\cxtitle{A determinant inequality for sums of Hermitian exponentials}
\begin{cxcredits}
\posedby{S.~Sra}
\foundby{GPT Astra (Pro)}{2026-10-30}
\formalizedby{GPT Astra (Pro)}
\auditedby{GPT Astra (Pro)}
\contributedby{S.~Sra}
\end{cxcredits}

\begin{cxcontext}
All matrices in a comparison have the same finite order. A Hermitian matrix
satisfies $H=H^*$, where $*$ denotes conjugate transpose, and $e^H$ is the
usual matrix exponential. The notation $\det(e^H+e^K)^2$ means the square
of the determinant, equivalently $\det((e^H+e^K)^2)$, not an entrywise square.
For Hermitian $H,K$, set
\[
 A=e^H,\qquad B=e^K,\qquad E=e^{H+K},\qquad D=B^{-1}=e^{-K}.
\]
These exponentials are positive definite. Products such as $DA$ need not
be Hermitian; the determinant identities used below do not require them to be.
\end{cxcontext}

\begin{cxsource}{hermitian-exponential-determinant}
S.~Sra, Conjecture~4.6, posed in 2015 while studying MathOverflow question
224121 \citep{Wolfgang2015MO}
\end{cxsource}

\begin{cxstatement}{hermitian-exponential-determinant}
For Hermitian matrices $H,K$ of the same order,
\[
 \det(e^H+e^K)^2
 \leq \det\bigl(e^{2H}+e^{2K}+2e^{H+K}\bigr).
\]
\end{cxstatement}

\begin{cxsummary}{hermitian-exponential-determinant}
A proposed determinant inequality comparing sums of Hermitian exponentials.
\end{cxsummary}

\begin{cxcertificate}{hermitian-exponential-determinant}
Integer $3\times3$ symmetric inputs; exact rational Taylor intervals certify normalized right side $<1007/1000<101/100<$ normalized left side.
\end{cxcertificate}

\begin{cxrefutation}
\begin{theorem}[\statusfalse: Hermitian exponential determinant comparison]\label{thm:hermitian-exponential-determinant}
The assertion in Problem~\ref{ctx:hermitian-exponential-determinant} is false
already for real symmetric $3\times3$ matrices.
\end{theorem}
\begin{proof}
Take the exact integer matrices
\[
 H=\begin{pmatrix}
 -26&10&-10\\
 10&-16&20\\
 -10&20&-26
 \end{pmatrix},
 \qquad
 K=\begin{pmatrix}
 0&0&0\\
 0&10&0\\
 0&0&20
 \end{pmatrix}.
\]
Both are real symmetric, hence Hermitian, of the same order. Define
\[
 L=\frac{\det(A+B)^2}{\det(B)^2},\qquad
 R=\frac{\det(A^2+B^2+2E)}{\det(B)^2}.
\]
Here $\det(B)^2=e^{2\operatorname{tr}K}=e^{60}>0$, so the proposed
inequality is equivalent to $L\leq R$. Left multiplication and congruence,
respectively, give the exact identities
\[
 L=\det(I+DA)^2,\qquad
 R=\det(I+DA^2D+2DED).
\]
The supplied verifier certifies
\begin{align*}
 1.0106259614312670597488454657
 &<L<1.0106259614312670597488454658,\\
 1.0066369661468304923537507020
 &<R<1.0066369661468304923537507021.
\end{align*}
Every terminating decimal in these strict bounds denotes an exact rational
number, not a rounded estimate. In particular,
\[
 R<\frac{1007}{1000}<\frac{101}{100}<L,
 \qquad L-R>\frac{3}{1000}.
\]
This reverses the claimed comparison.

For completeness, here is the error control behind the certificate.
For a rational matrix $X$ put
\[
 r=\|X\|_\infty=\max_i\sum_j|X_{ij}|,
 \qquad T_m(X)=\sum_{j=0}^{m}\frac{X^j}{j!}.
\]
Submultiplicativity and the power-series definition imply
\[
 \|e^X-T_m(X)\|_\infty
 \leq\sum_{j=m+1}^{\infty}\frac{r^j}{j!}
 \leq\frac{r^{m+1}}{(m+1)!}\,
       \frac{1}{1-r/(m+2)}
 =:\varepsilon_m(r),\qquad r<m+2.
\]
The last bound is a geometric-series bound: after the first omitted term,
each successive ratio is at most $r/(m+2)$. Therefore every entry of $e^X$
lies in the rational interval
$[(T_m(X))_{ij}-\varepsilon_m(r),
  (T_m(X))_{ij}+\varepsilon_m(r)]$.

The verifier takes $m=256$. The required norms are
$\|H\|_\infty=56$, $\|H+K\|_\infty=46$, and $0,10,20$ for the three
scalar exponents forming $D$. Each is strictly below $258$.
The off-diagonal entries of $D$ are exactly zero. Rational interval addition,
subtraction, and multiplication propagate these enclosures through the two
normalized determinant expressions, using the six-term determinant formula
in order three. Before squaring the left determinant interval, its strict
positivity is certified. Interval endpoints and every inequality test use
integers and \texttt{fractions.Fraction}; there is no unbounded truncation
error and no machine-precision comparison. All displayed enclosing bounds
and the strict separation above are checked, not merely printed.
\end{proof}

\begin{remark}[The linear coefficient and a log-determinant normalization]
The following coefficientwise strengthening is recorded only as an auxiliary
observation, not as a second archival result.
For fixed $A,B,E$, let
\[
 \Delta(t)=\det(B^2+2tE+t^2A^2)-\det(B+tA)^2.
\]
Jacobi's determinant differentiation formula at $t=0$ gives
\[
 \frac{[t]\Delta(t)}{\det(B)^2}
 =2\bigl(\operatorname{tr}(B^{-2}E)
          -\operatorname{tr}(B^{-1}A)\bigr)
 =2\bigl(\operatorname{tr}(DED)-\operatorname{tr}(DA)\bigr).
\]
For the same exact witness, the verifier certifies that this number $c_1$
satisfies
\[
 -0.0081120478335667666163258676
 <c_1<
 -0.0081120478335667666163258675
 <-\frac1{125}<0.
\]
Thus even the linear coefficient has the wrong sign.

This is also the obstruction exposed by a log-determinant normalization.
Put $S_t=B^2+t^2A^2$. Since
$(B+tA)^2=S_t+t(AB+BA)$, cancellation of $\det S_t$ gives, for $t$ near zero,
\begin{align*}
 F(t)
 &=\operatorname{tr}\log\bigl(I+2tS_t^{-1/2}ES_t^{-1/2}\bigr)\\
 &\quad-\operatorname{tr}\log\bigl(I+tS_t^{-1/2}(AB+BA)S_t^{-1/2}\bigr)\\
 &=c_1t+O(t^2).
\end{align*}
The logarithms are defined in a neighborhood of zero by positivity and
continuity. The expansion is used to explain the discovery, not as a
substitute for the exact certificate at $t=1$.
\end{remark}

\begin{remark}[An analytic family explaining the failure]
A separate limiting construction explains why counterexamples exist.
Set
\[
 M=\begin{pmatrix}0&1&-1\\1&1&2\\-1&2&0\end{pmatrix},
 \qquad N=\operatorname{diag}(0,1,2),
 \qquad \lambda=\lambda_{\max}(M).
\]
The unit vector $v=(0,1,1)^T/\sqrt2$ has $v^TMv=5/2$, so
$\lambda\geq5/2>2$. No eigenvector of $M$ can have first coordinate zero: writing its eigenvalue
as $\mu$, the first row would give $u_2=u_3$, and the second and third rows
would then give simultaneously $\mu u_2=3u_2$ and $\mu u_2=2u_2$.
This would force the entire vector to vanish. Consequently each eigenspace
is one-dimensional (otherwise a linear combination would have first
coordinate zero). Let $P=uu^T$ be the rank-one top spectral projection,
with $\|u\|=1$, and put $Q=e_1e_1^T$, $\alpha=|u_1|^2>0$.

For $r>0$ define
\[
 H_r=r(M-\lambda I),\quad K_r=rN,\quad
 A_r=e^{H_r},\quad B_r=e^{K_r},\quad
 D_r=B_r^{-1},\quad E_r=e^{H_r+K_r}.
\]
Then $A_r\longrightarrow P$ and $D_r\longrightarrow Q$. The matrix
$T=M+N$ has characteristic polynomial $(x-4)(x^2-2)$ and top unit
eigenvector $v$. Since this vector has zero first coordinate,
\begin{align*}
 \|D_rE_rD_r\|_2
 &\leq e^{(\sqrt2-\lambda)r}
       +e^{(4-\lambda)r}\|D_rv\|^2\\
 &=e^{(\sqrt2-\lambda)r}
   +\tfrac12 e^{(2-\lambda)r}
   +\tfrac12 e^{-\lambda r}
 \longrightarrow0.
\end{align*}
Here $\|\cdot\|_2$ is the operator norm. The contribution from the two
remaining eigenspaces is bounded by $e^{(\sqrt2-\lambda)r}$, since
$\|D_r\|_2=1$.

It follows that, coefficientwise as polynomials in $t$,
\begin{align*}
 \frac{\det(B_r^2+2tE_r+t^2A_r^2)}{\det(B_r)^2}
 &=\det(I+2tD_rE_rD_r+t^2D_rA_r^2D_r)
 \longrightarrow1+\alpha t^2,\\
 \frac{\det(B_r+tA_r)^2}{\det(B_r)^2}
 &=\det(I+tD_rA_r)^2
 \longrightarrow(1+\alpha t)^2.
\end{align*}
Coefficientwise convergence follows because these coefficients are polynomial
functions of the converging matrix entries. The normalized difference tends to
\[
 -2\alpha t+\alpha(1-\alpha)t^2.
\]
At $t=1$ its limit is $-\alpha(1+\alpha)<0$, so the original inequality
fails for all sufficiently large $r$. Its linear coefficient tends to
$-2\alpha<0$. This analytic explanation is not needed to validate the fixed
integer witness above; in particular, no numerical value of $\lambda$ enters
the executable certificate.
\end{remark}

\begin{remark}[Scope]
The certified conclusion concerns the displayed universal assertion for
same-order Hermitian matrices, with the direction $\leq$ and with the term
$2e^{H+K}$ unchanged.

Commuting pairs satisfy equality, and this case does not dispute that fact.
No claim is made that dimension three is minimal, nor is the dimension-two
case resolved here. The negative coefficient above is auxiliary, not a
separately provenanced claim. No novelty or priority is asserted.
\end{remark}
\end{cxrefutation}
